\begin{table}[b]
    \centering
    \begin{ruledtabular}
    \begin{tabular}{cccccc}
        $L$ & $M$ & BC-Z & $\xi_{n=1}(T=0.7)$ & $\Delta\xi_{n=1}(T=0.7)$ & $\xi_{n=1}(T\simeq T_\mathrm{c})$ \\\hline
         $4$ & $192$ & PBC & 8.316(15) & 4.402(18) & 3.914(16)\\
         $6$ & $192$ & PBC & 13.363(24) & 7.492(22) & 5.871(12)\\
         $8$ & $192$ & PBC & 18.87(3) & 11.06(3) & 7.807(13)\\
         $12$ & $320$ & PBC & 30.79(6) & 19.04(6) & 11.742(23)\\
         $16$ & $512$ & PBC & 44.6(4) & 28.7(4) & 15.96(14)\\
         $16$ & $48$ & OBC & 44.3(4) & 28.6(4) & 15.70(10)\\
         $24$ & $88$ & OBC & 73.0(16) & 49.1(16) & 23.9(4)\\
   \end{tabular}
    \end{ruledtabular}
    \caption{$\xi_{n=1}$ and $\Delta\xi_{n=1}$ for the different system sizes, $L^2\times M$, boundary conditions along the Z-axis (BC-Z), and temperatures. Data is represented in Fig.~\ref{fig:xi_vs_L}.}
    \label{tab:xi}
\end{table}
